% GENERATED FILE. Edit canonical lessons, then run npm run generate:books.
% canonical-source-hash: fnv1a64:aa37f0d662c84bae
% canonical-lessons: RU-C98-frying-pan, RU-C98-bucket, RU-C98-rope, RU-C98-nail, RU-C98-hammer

\chapter{Frying pan, Bucket, Rope, Nail, Hammer}
\label{ch:ru-frying-pan-bucket-rope-nail-hammer}
\hlchaptermodality{98}

\begin{chapteropening}
\textbf{By the end of this chapter:} \emph{I can say a frying pan, a bucket, a rope, a nail and a hammer.}

Complete the last lesson of chapter 98: I can say a frying pan, a bucket, a rope, a nail and a hammer.
\end{chapteropening}

\section[\texorpdfstring{\ru{сковорода} (skovorodá)}{сковорода (skovorodá)}]{\ru{сковорода} (skovorodá) — a frying pan}

\label{lesson:RU-C98-frying-pan}



\begin{quote}
Before the new one: say the Russian for a wardrobe, then the Russian for a saucepan.
\end{quote}

\subsection*{You'll want to know: \ru{сковорода}}
\textbf{\ru{сковорода}} (\emph{skovorodá}) — \textquotedblleft{}a frying pan\textquotedblright{}.

\begin{cousinweb}
\textbf{\ru{сковорода}} $\leftarrow$ Proto-Slavic \emph{*skovorda} — inherited, with Church Slavonic \emph{\ru{сковрада}}. Vasmer calls its further origin unclear; a comparison with Old High German \emph{scart-îsan}, \textquotedblleft{}a frying pan\textquotedblright{}, has been made, and Szemerényi derived it from an Iranian word for \textquotedblleft{}bowl\textquotedblright{}.
\end{cousinweb}

\begin{grammarlens}[title={What you've built}]
One more word for this chapter.
\end{grammarlens}

\subsection*{Your turn}
\begin{itemize}
\raggedright
  \item \emph{Say it:} \emph{skovorodá}
  \item \emph{Say it:} \emph{skovorodá}, once more
  \item \emph{Say it:} say \emph{skovorodá}
  \item \emph{Recall:} say the Russian for a wardrobe, then the Russian for a saucepan, then say \emph{skovorodá} again
\end{itemize}

\begin{tcolorbox}[breakable,colback=teal!4,colframe=teal!35!black,title={Before you move on}]
What does \ru{сковорода} mean? (\textquotedblleft{}A frying pan\textquotedblright{}.) Say it once more.
\end{tcolorbox}
\section[\texorpdfstring{\ru{ведро} (vedró)}{ведро (vedró)}]{\ru{ведро} (vedró) — a bucket}

\label{lesson:RU-C98-bucket}



\begin{quote}
Before the new one: say the Russian for a saucepan, then the Russian for a frying pan.
\end{quote}

\subsection*{You'll want to know: \ru{ведро}}
\textbf{\ru{ведро}} (\emph{vedró}) — \textquotedblleft{}a bucket\textquotedblright{}.

\begin{cousinweb}
\textbf{\ru{ведро}} $\leftarrow$ Proto-Slavic \emph{*vědro} — inherited. Vasmer relates it to \emph{\ru{вода}}, \textquotedblleft{}water\textquotedblright{}, German \emph{Wasser} and Greek \emph{hydría}, \textquotedblleft{}a water jar\textquotedblright{}: the bucket as the water vessel.
\end{cousinweb}

\begin{grammarlens}[title={What you've built}]
One more word for this chapter.
\end{grammarlens}

\subsection*{Your turn}
\begin{itemize}
\raggedright
  \item \emph{Say it:} \emph{vedró}
  \item \emph{Say it:} \emph{vedró}, once more
  \item \emph{Say it:} say \emph{vedró}
  \item \emph{Recall:} say the Russian for a saucepan, then the Russian for a frying pan, then say \emph{vedró} again
\end{itemize}

\begin{tcolorbox}[breakable,colback=teal!4,colframe=teal!35!black,title={Before you move on}]
What does \ru{ведро} mean? (\textquotedblleft{}A bucket\textquotedblright{}.) Say it once more.
\end{tcolorbox}
\section[\texorpdfstring{\ru{верёвка} (veryóvka)}{верёвка (veryóvka)}]{\ru{верёвка} (veryóvka) — a rope}

\label{lesson:RU-C98-rope}



\begin{quote}
Before the new one: say the Russian for a frying pan, then the Russian for a bucket.
\end{quote}

\subsection*{You'll want to know: \ru{верёвка}}
\textbf{\ru{верёвка}} (\emph{veryóvka}) — \textquotedblleft{}a rope\textquotedblright{}.

\begin{cousinweb}
\textbf{\ru{верёвка}} is built on Old Russian \emph{\ru{вьрвь}}, \textquotedblleft{}a cord\textquotedblright{}, with \emph{-\ru{ка}}. Vasmer relates it to Lithuanian \emph{virvė}, \textquotedblleft{}rope\textquotedblright{}, and \emph{vérti}, \textquotedblleft{}to thread\textquotedblright{}, and Greek \emph{erýō}, \textquotedblleft{}I drag\textquotedblright{}.
\end{cousinweb}

\begin{grammarlens}[title={What you've built}]
One more word for this chapter.
\end{grammarlens}

\subsection*{Your turn}
\begin{itemize}
\raggedright
  \item \emph{Say it:} \emph{veryóvka}
  \item \emph{Say it:} \emph{veryóvka}, once more
  \item \emph{Say it:} say \emph{veryóvka}
  \item \emph{Recall:} say the Russian for a frying pan, then the Russian for a bucket, then say \emph{veryóvka} again
\end{itemize}

\begin{tcolorbox}[breakable,colback=teal!4,colframe=teal!35!black,title={Before you move on}]
What does \ru{верёвка} mean? (\textquotedblleft{}A rope\textquotedblright{}.) Say it once more.
\end{tcolorbox}
\section[\texorpdfstring{\ru{гвоздь} (gvozd)}{гвоздь (gvozd)}]{\ru{гвоздь} (gvozd) — a nail}

\label{lesson:RU-C98-nail}



\begin{quote}
Before the new one: say the Russian for a bucket, then the Russian for a rope.
\end{quote}

\subsection*{You'll want to know: \ru{гвоздь}}
\textbf{\ru{гвоздь}} (\emph{gvozd}) — \textquotedblleft{}a nail\textquotedblright{}.

\begin{cousinweb}
\textbf{\ru{гвоздь}} $\leftarrow$ Proto-Slavic \emph{*gvozdĭ} — inherited. Vasmer will not separate it from old Slavic \emph{gvozd}, \textquotedblleft{}a forest\textquotedblright{}, and compares German \emph{Quaste}, \textquotedblleft{}a tassel\textquotedblright{}, and Old Norse \emph{kvistr}, \textquotedblleft{}a twig\textquotedblright{}: a nail was first a wooden peg.
\end{cousinweb}

\begin{grammarlens}[title={What you've built}]
One more word for this chapter.
\end{grammarlens}

\subsection*{Your turn}
\begin{itemize}
\raggedright
  \item \emph{Say it:} \emph{gvozd}
  \item \emph{Say it:} \emph{gvozd}, once more
  \item \emph{Say it:} say \emph{gvozd}
  \item \emph{Recall:} say the Russian for a bucket, then the Russian for a rope, then say \emph{gvozd} again
\end{itemize}

\begin{tcolorbox}[breakable,colback=teal!4,colframe=teal!35!black,title={Before you move on}]
What does \ru{гвоздь} mean? (\textquotedblleft{}A nail\textquotedblright{}.) Say it once more.
\end{tcolorbox}
\section[\texorpdfstring{\ru{молоток} (molotók)}{молоток (molotók)}]{\ru{молоток} (molotók) — a hammer}

\label{lesson:RU-C98-hammer}



\begin{quote}
Before the new one: say the Russian for a rope, then the Russian for a nail.
\end{quote}

\subsection*{You'll want to know: \ru{молоток}}
\textbf{\ru{молоток}} (\emph{molotók}) — \textquotedblleft{}a hammer\textquotedblright{}.

\begin{cousinweb}
\textbf{\ru{молоток}} is \emph{\ru{молот}}, \textquotedblleft{}a sledgehammer\textquotedblright{}, with the diminutive \emph{-\ru{ок}}. Vasmer thinks \emph{\ru{молот}} probably comes from \emph{\ru{молоть}}, \textquotedblleft{}to grind\textquotedblright{}, so \textquotedblleft{}the crusher\textquotedblright{}, and compares Latin \emph{malleus}, \textquotedblleft{}a hammer\textquotedblright{}.
\end{cousinweb}

\begin{grammarlens}[title={What you've built}]
One more word for this chapter.
\end{grammarlens}

\subsection*{Your turn}
\begin{itemize}
\raggedright
  \item \emph{Say it:} \emph{molotók}
  \item \emph{Say it:} \emph{molotók}, once more
  \item \emph{Say it:} say \emph{molotók}
  \item \emph{Recall:} say the Russian for a rope, then the Russian for a nail, then say \emph{molotók} again
\end{itemize}

\begin{tcolorbox}[breakable,colback=teal!4,colframe=teal!35!black,title={Before you move on}]
What does \ru{молоток} mean? (\textquotedblleft{}A hammer\textquotedblright{}.) Say it once more.
\end{tcolorbox}
